% ==============================================================================
\section{Одномерное дискретное преобразование Фурье (ДПФ)}
\label{sec:1d_fourier}
% ==============================================================================

Дискретное преобразование Фурье является фундаментальным инструментом анализа сигналов, позволяющим осуществлять переход от временной (или пространственной) области к частотной. 

\begin{tcbdefinition}[Матрица Фурье]
\label{def:fourier_matrix}
    Пусть $\omega_n = e^{-\frac{2\pi i}{n}}$ — первообразный корень степени $n$ из единицы. \textit{Матрицей Фурье} порядка $n$ называется матрица $\bm{F}_n \in \C^{n \times n}$, элементы которой задаются соотношением:
    \begin{equation}
        (\bm{F}_n)_{p,q} = \omega_n^{p q} = e^{-\frac{2\pi i}{n} p q}, \quad p, q \in \{0, 1, \dots, n-1\}.
        \label{eq:fourier_elements}
    \end{equation}
    В развернутом виде матрица принимает форму матрицы Вандермонда:
    \begin{equation}
        \bm{F}_n = 
        \begin{pNiceMatrix}
            1 & 1 & 1 & \Cdots & 1 \\
            1 & \omega_n & \omega_n^2 & \Cdots & \omega_n^{n-1} \\
            1 & \omega_n^2 & \omega_n^4 & \Cdots & \omega_n^{2(n-1)} \\
            \Vdots & \Vdots & \Vdots & \Ddots & \Vdots \\
            1 & \omega_n^{n-1} & \omega_n^{2(n-1)} & \Cdots & \omega_n^{(n-1)(n-1)}
        \end{pNiceMatrix}.
        \label{eq:fourier_matrix}
    \end{equation}
\end{tcbdefinition}

\begin{tcbtheorem}[Свойства матрицы Фурье]
\label{thm:fourier_props}
    Матрица Фурье $\bm{F}_n$ обладает следующими алгебраическими и геометрическими свойствами:
    \begin{itemize}
        \item \textbf{Симметричность:} $\bm{F}_n = \bm{F}_n\T$.
        \item \textbf{Почти унитарность:} $\bm{F}_n\H \bm{F}_n = n \bm{I}$, следовательно, матрица $\frac{1}{\sqrt{n}} \bm{F}_n$ является унитарной, а обратная матрица вычисляется как $\bm{F}_n^{-1} = \frac{1}{n} \bm{F}_n\H$.
        \item \textbf{Сохранение углов:} $\inner{\bm{F}_n \bm{x}, \bm{F}_n \bm{y}} = n \inner{\bm{x}, \bm{y}}$ для любых $\bm{x}, \bm{y} \in \C^n$.
        \item \textbf{Комплексная сопряженность столбцов:} Пусть $\bm{f}_q$ — $q$-й столбец матрицы $\bm{F}_n$. Тогда $\overline{\bm{f}_q} = \bm{f}_{n-q}$ для $q \in \{1, \dots, n-1\}$.
    \end{itemize}
\end{tcbtheorem}

\begin{tcbproof}[теоремы~\ref{thm:fourier_props}]
    \textit{Симметричность} тривиально следует из определения~\cref{eq:fourier_elements}, так как $(\bm{F}_n)_{p,q} = \omega_n^{pq} = \omega_n^{qp} = (\bm{F}_n)_{q,p}$.
    
    Для доказательства \textit{почти унитарности} рассмотрим скалярное произведение $p$-го и $q$-го столбцов:
    \begin{equation}
        \bm{f}_p\H \bm{f}_q = \sum_{j=0}^{n-1} \overline{(\bm{F}_n)_{j,p}} (\bm{F}_n)_{j,q} = \sum_{j=0}^{n-1} e^{\frac{2\pi i}{n} j p} e^{-\frac{2\pi i}{n} j q} = \sum_{j=0}^{n-1} e^{\frac{2\pi i}{n} j (p-q)}.
    \end{equation}
    При $p = q$ показатель экспоненты равен нулю, и сумма тривиально равна $n$. При $p \neq q$ имеем сумму геометрической прогрессии:
    \begin{equation}
        \sum_{j=0}^{n-1} \left( e^{\frac{2\pi i}{n} (p-q)} \right)^j = \frac{1 - e^{2\pi i (p-q)}}{1 - e^{\frac{2\pi i}{n} (p-q)}} = 0,
    \end{equation}
    поскольку числитель равен нулю ($e^{2\pi i k} = 1$ для целых $k$), а знаменатель отличен от нуля. Таким образом, $\bm{F}_n\H \bm{F}_n = n \bm{I}$.
    
    Свойство \textit{сохранения углов} вытекает из предыдущего:
    \begin{equation}
        \inner{\bm{F}_n \bm{x}, \bm{F}_n \bm{y}} = (\bm{F}_n \bm{y})\H (\bm{F}_n \bm{x}) = \bm{y}\H \bm{F}_n\H \bm{F}_n \bm{x} = n \bm{y}\H \bm{x} = n \inner{\bm{x}, \bm{y}}.
    \end{equation}
    
    Для доказательства \textit{сопряженности столбцов} рассмотрим $k$-й элемент вектора $\bm{f}_{n-q}$:
    \begin{equation}
        (\bm{f}_{n-q})_k = \exp\left(-\frac{2\pi i}{n} k(n-q)\right) = \exp(-2\pi i k) \exp\left(\frac{2\pi i}{n} k q\right) = 1 \cdot \overline{(\bm{f}_q)_k},
    \end{equation}
    что завершает доказательство.
\end{tcbproof}

\begin{tcbinsight}[Смысл базиса Фурье]
    Матрицу $\bm{F}_n^{-1}$ целесообразно рассматривать как матрицу перехода к новому базису. Координаты базисных векторов (столбцов $\bm{F}_n^{-1}$) образованы дискретными отсчетами гармонических функций (косинусов и синусов). Свойство сопряженности столбцов из \cref{thm:fourier_props} означает, что при применении ДПФ к вещественному сигналу $\bm{x} \in \R^n$, выходной вектор $\bm{X} = \bm{F}_n \bm{x}$ обладает симметрией $\bm{X}_k = \overline{\bm{X}_{n-k}}$.
\end{tcbinsight}

% ==============================================================================
\section{Структурированные матрицы и быстрые алгоритмы}
\label{sec:structured_matrices}
% ==============================================================================

Быстрое преобразование Фурье (БПФ) снижает сложность вычисления матрично-векторного произведения $\bm{F}_n \bm{x}$ с $O(n^2)$ до $O(n \log n)$. Это свойство радикально ускоряет работу со структурированными матрицами, такими как циркулянты и матрицы Тёплица.

% ==============================================================================
\subsection{Циркулянтные матрицы}
\label{sec:circulant}
% ==============================================================================

\begin{tcbdefinition}[Циркулянт]
\label{def:circulant}
    Матрица $\bm{C} \in \C^{n \times n}$ называется \textit{циркулянтом}, если каждая её последующая строка получается из предыдущей циклическим сдвигом вправо на одну позицию:
    \begin{equation}
        \bm{C} = 
        \begin{pNiceMatrix}
            c_0 & c_{n-1} & \Cdots & c_1 \\
            c_1 & c_0 & \Cdots & c_2 \\
            \Vdots & \Vdots & \Ddots & \Vdots \\
            c_{n-1} & c_{n-2} & \Cdots & c_0
        \end{pNiceMatrix}.
    \end{equation}
\end{tcbdefinition}

\begin{tcbtheorem}[Диагонализуемость циркулянтов]
\label{thm:circulant_diag}
    Любая циркулянтная матрица $\bm{C}$ диагонализуется матрицей Фурье:
    \begin{equation}
        \bm{C} = \bm{F}_n^{-1} \diag(\bm{F}_n \bm{c}) \bm{F}_n,
    \end{equation}
    где $\bm{c} = (c_0, c_1, \dots, c_{n-1})\T$ — первый столбец матрицы $\bm{C}$. Собственные значения матрицы $\bm{C}$ суть компоненты вектора $\bm{F}_n \bm{c}$.
\end{tcbtheorem}

\begin{tcbproof}[теоремы~\ref{thm:circulant_diag}]
    Непосредственная проверка показывает, что применение $\bm{C}$ к столбцам $\bm{F}_n^{-1}$ равносильно умножению этих столбцов на скаляры, представляющие собой элементы ДПФ от первого столбца $\bm{c}$. 
\end{tcbproof}

\begin{tcbremark}
    Из \cref{thm:circulant_diag} следует коммутативность алгебры циркулянтов: $\bm{C}_1 \bm{C}_2 = \bm{C}_2 \bm{C}_1$. Необратимость циркулянта эквивалентна наличию хотя бы одного нуля в спектре $\diag(\bm{F}_n \bm{c})$.
\end{tcbremark}

\begin{tcbalgorithmbox}[Быстрое умножение циркулянта на вектор]
\label{alg:circulant_mult}
    \textbf{Вход:} Вектор $\bm{c}$ (первый столбец циркулянта $\bm{C}$), вектор $\bm{x} \in \C^n$.\\
    \textbf{Выход:} Вектор $\bm{y} = \bm{C}\bm{x}$.
    \begin{enumerate}
        \item Вычислить БПФ вектора $\bm{c}$: $\bm{\lambda} \gets \bm{F}_n \bm{c}$ (затраты $O(n \log n)$).
        \item Вычислить БПФ вектора $\bm{x}$: $\bm{\hat{x}} \gets \bm{F}_n \bm{x}$ (затраты $O(n \log n)$).
        \item Выполнить поэлементное умножение: $\bm{\hat{y}}_i \gets \bm{\lambda}_i \cdot \bm{\hat{x}}_i$ (затраты $O(n)$).
        \item Выполнить обратное БПФ: $\bm{y} \gets \bm{F}_n^{-1} \bm{\hat{y}}$ (затраты $O(n \log n)$).
    \end{enumerate}
    Итоговая асимптотическая сложность: $O(n \log n)$.
\end{tcbalgorithmbox}

% ==============================================================================
\subsection{Тёплицевы матрицы}
\label{sec:toeplitz}
% ==============================================================================

\begin{tcbdefinition}[Матрица Тёплица]
\label{def:toeplitz}
    Матрица $\bm{T} \in \C^{n \times n}$ называется \textit{Тёплицевой}, если элементы на её диагоналях постоянны ($T_{i,j} = t_{i-j}$):
    \begin{equation}
        \bm{T} = 
        \begin{pNiceMatrix}
            t_0 & t_{-1} & \Cdots & t_{-(n-1)} \\
            t_1 & t_0 & \Ddots & \Vdots \\
            \Vdots & \Ddots & \Ddots & t_{-1} \\
            t_{n-1} & \Cdots & t_1 & t_0
        \end{pNiceMatrix}.
    \end{equation}
\end{tcbdefinition}

Произведение Тёплицевых матриц не обязательно сохраняет структуру. Однако умножение на вектор можно выполнить эффективно.

\begin{tcbalgorithmbox}[Умножение Тёплицевой матрицы на вектор]
\label{alg:toeplitz_mult}
    \textbf{Вход:} Тёплицева матрица $\bm{T} \in \C^{n \times n}$, вектор $\bm{x} \in \C^n$.\\
    \textbf{Выход:} Вектор $\bm{y} = \bm{T}\bm{x}$.
    \begin{enumerate}
        \item Сформировать циркулянтную матрицу $\bm{\hat{C}} \in \C^{2n \times 2n}$, вложив в неё матрицу $\bm{T}$. Первый столбец $\bm{\hat{c}}$ матрицы $\bm{\hat{C}}$ имеет вид:
        \begin{equation*}
            \bm{\hat{c}} = \begin{bmatrix} t_0 & t_1 & \dots & t_{n-1} & 0 & t_{-(n-1)} & \dots & t_{-1} \end{bmatrix}\T.
        \end{equation*}
        \item Дополнить вектор $\bm{x}$ нулями до размера $2n$: 
        $\bm{\hat{x}} = \begin{bmatrix} \bm{x}\T & 0 & \dots & 0 \end{bmatrix}\T$.
        \item Вычислить произведение $\bm{\hat{y}} = \bm{\hat{C}} \bm{\hat{x}}$, используя \cref{alg:circulant_mult}.
        \item Извлечь первые $n$ элементов из $\bm{\hat{y}}$: $\bm{y} = \bm{\hat{y}}_{1 \dots n}$.
    \end{enumerate}
\end{tcbalgorithmbox}

\begin{tcbexample}[Свертка сигналов]
\label{ex:convolution}
    Вычисление линейной свертки дискретных сигналов эквивалентно умножению Тёплицевой матрицы (построенной из отсчетов одного сигнала) на вектор-столбец второго сигнала. Применение \cref{alg:toeplitz_mult} позволяет вычислять такие свертки за $O(n \log n)$.
\end{tcbexample}

% ==============================================================================
\section{Быстрое умножение многочленов}
\label{sec:polynomials}
% ==============================================================================

Пусть заданы два многочлена $f(x)$ и $g(x)$ степени $n$. Их коэффициенты представимы в виде векторов $\bm{a}, \bm{b} \in \C^{n+1}$. Прямое вычисление коэффициентов произведения $h(x) = f(x)g(x)$ требует $O(n^2)$ операций. Использование ДПФ позволяет решить задачу асимптотически быстрее.

\begin{tcbalgorithmbox}[Быстрое умножение многочленов]
\label{alg:poly_mult}
    \textbf{Вход:} Коэффициенты многочленов $\bm{a}, \bm{b} \in \C^{n+1}$.\\
    \textbf{Выход:} Вектор коэффициентов $\bm{c} \in \C^{2n+1}$ многочлена $f(x)g(x)$.
    \begin{enumerate}
        \item Дополнить векторы $\bm{a}$ и $\bm{b}$ нулями до размера $N \geq 2n+1$ (обычно выбирается ближайшая степень двойки).
        \item Вычислить значения многочленов в корнях из единицы степени $N$ с помощью БПФ: 
              $\bm{\hat{f}} \gets \bm{F}_N \bm{a}$, $\bm{\hat{g}} \gets \bm{F}_N \bm{b}$.
        \item Выполнить поэлементное умножение: $\bm{\hat{h}}_i \gets \bm{\hat{f}}_i \cdot \bm{\hat{g}}_i$ для $i = 0, \dots, N-1$. (Это дает значения многочлена $h(x)$ в узлах интерполяции).
        \item Найти коэффициенты произведения интерполяцией с помощью обратного БПФ: 
              $\bm{c} \gets \bm{F}_N^{-1} \bm{\hat{h}}$.
    \end{enumerate}
\end{tcbalgorithmbox}

% ==============================================================================
\section{Применение ДПФ: Фильтрация сигналов}
\label{sec:signal_filtering}
% ==============================================================================

\begin{tcbinsight}[Частотная фильтрация]
    Сигналы (например, данные о температуре или звуковые волны) содержат периодические паттерны, смешанные с шумом. Энергия (норма $\norm{\bm{x}}_2^2$) сигнала сохраняется при ДПФ с точностью до множителя (теорема Парсеваля). Если основная "масса" сигнала сосредоточена в узких пиках Фурье-образа, а шум равномерно распределен по спектру, то обнуление спектра вне пиков восстанавливает полезный сигнал.
\end{tcbinsight}

Для фильтрации можно применять как нелинейные (пороговые), так и линейные методы.

\begin{tcbalgorithmbox}[Пороговая и взвешенная фильтрация одномерного сигнала]
\label{alg:signal_filtering}
    \textbf{Вход:} Зашумленный сигнал $\bm{x} \in \R^n$.\\
    \textbf{Выход:} Очищенный сигнал $\bm{\tilde{x}} \in \R^n$.
    \begin{enumerate}
        \item Переход в частотную область: $\bm{X} \gets \bm{F}_n \bm{x}$.
        \item Применение фильтра:
        \begin{itemize}
            \item \textit{Пороговый (нелинейный) фильтр:} $\bm{\tilde{X}}_k \gets 0$, если $\abs{\bm{X}_k} < \tau$.
            \item \textit{Сглаживающий (линейный) фильтр:} умножение на окно, например, Гауссову "шапочку": $\bm{\tilde{X}}_k \gets \bm{X}_k \cdot \exp\left(-\alpha \frac{(\min(k, n-k))^2}{2}\right)$, что сохраняет низкие частоты и подавляет высокие (учитывая симметрию спектра).
        \end{itemize}
        \item Возврат во временную область: $\bm{\tilde{x}} \gets \operatorname{Re}(\bm{F}_n^{-1} \bm{\tilde{X}})$.
    \end{enumerate}
\end{tcbalgorithmbox}

% ==============================================================================
\section{Двумерное преобразование Фурье и обработка изображений}
\label{sec:2d_fourier}
% ==============================================================================

Двумерные объекты (изображения) анализируются посредством двумерного ДПФ, которое эквивалентно последовательному применению одномерного ДПФ к строкам, а затем к столбцам матрицы изображения.

\begin{tcbdefinition}[Двумерное дискретное преобразование Фурье]
\label{def:2d_fourier}
    Пусть $\bm{X} \in \C^{m \times n}$ — матрица изображения. Её двумерное ДПФ $\bm{Y} \in \C^{m \times n}$ определяется как:
    \begin{equation}
        \bm{Y} = \bm{F}_m \bm{X} \bm{F}_n\T.
    \end{equation}
    В терминах кронекеровского произведения и векторизации это записывается как:
    \begin{equation}
        \operatorname{vec}(\bm{Y}) = (\bm{F}_n \otimes \bm{F}_m) \operatorname{vec}(\bm{X}).
    \end{equation}
    Поэлементно преобразование имеет вид:
    \begin{equation}
        Y_{p_1, p_2} = \sum_{q_1=0}^{m-1} \sum_{q_2=0}^{n-1} \exp\left(-\frac{2\pi i p_1 q_1}{m} - \frac{2\pi i p_2 q_2}{n}\right) X_{q_1, q_2}.
    \end{equation}
\end{tcbdefinition}

\begin{tcbremark}
    Базисные функции двумерного ДПФ представляют собой двумерные комплексные экспоненты. Их вещественные части визуально выглядят как плоские волны (муаровые узоры). Ориентация и частота этих волн зависят от индексов $(p_1, p_2)$.
\end{tcbremark}

\begin{tcbexample}[Очистка текстуры холста]
\label{ex:canvas_removal}
    При оцифровке картин часто возникает паразитная периодическая структура — текстура холста. В частотной области (при визуализации логарифма модуля $\log_2(\abs{\bm{Y}} + 1)$) эта текстура проявляется в виде ярко выраженных "звездочек" или пиков вдали от центра (низких частот). Искусственное обнуление амплитуды в этих локальных окрестностях спектра с последующим обратным 2D ДПФ позволяет эффективно подавить текстуру холста, не размывая границы объектов на изображении.
\end{tcbexample}

\begin{tcbinsight}[Роль амплитуды и фазы]
    В матрице $\bm{Y} = \abs{\bm{Y}} \circ \exp(i\bm{\Phi})$ (где $\circ$ — поэлементное произведение) модуль $\abs{\bm{Y}}$ отвечает за амплитуду частот, а матрица фаз $\bm{\Phi} = \operatorname{arg}(\bm{Y})$ содержит информацию о пространственной локализации этих частот. Эксперименты с перекрестным обменом фазовых спектров двух разных изображений показывают, что результат обратного ДПФ визуально распознается как то изображение, \textit{фаза} которого была использована, несмотря на то что амплитуды заимствованы из совершенно другого сигнала.
\end{tcbinsight}
